\subsubsection{Real Sequences and Completeness Methods}

Real sequences are the first analytic testing device for limits, order, and completeness. This subsection develops convergence of real, complex, and Euclidean sequences; subsequences; Cauchy criteria on the real line; monotone convergence; upper and lower limits; and the Bolzano--Weierstrass method.

\begin{definition} \label{def:real-complex-sequence-convergence} \hlt{Convergent Real and Complex Sequences}\\
A real or complex sequence $(s_n)$ \hlt{converges} to $s$ if, for every $\varepsilon>0$, there is $N\in\N$ such that $n\geq N$ implies $\abs{s_n-s}<\varepsilon$.\\
Notation is $s_n\to s$ or $\lim\limits_{n\to\infty}s_n=s$. If no such $s$ exists, the sequence \hlt{diverges}.
\end{definition}

\begin{definition} \label{def:real-sequence-range-and-boundedness} \hlt{Range and Boundedness of a Real Sequence}\\
The \hlt{range} of $(s_n)$ is the set $\{s_n:n\in\N\}$. A real sequence is \hlt{bounded} if its range is bounded above and below in $\R$.
\end{definition}

\begin{theorem} \label{thm:complex-sequence-limit-algebra} \hlt{Limit Algebra for Complex Sequences}\\
Suppose $(s_n)$ and $(t_n)$ are complex sequences with $s_n\to s$ and $t_n\to t$. Then, for every $c\in\C$,
\begin{enumerate}[label=(\roman*)]
\setlength{\itemsep}{0pt}
\item $s_n+t_n\to s+t$;
\item $cs_n\to cs$ and $c+s_n\to c+s$;
\item $s_nt_n\to st$;
\item if $s\neq0$ and $s_n\neq0$ for all $n$, then $1/s_n\to1/s$.
\end{enumerate}
\begin{proof}
\begin{enumerate}[label=(\roman*)]
\setlength{\itemsep}{0pt}
\item Given $\varepsilon>0$, choose $N$ such that $\abs{s_n-s}<\varepsilon/2$, $\abs{t_n-t}<\varepsilon/2$ for $n\geq N$.\\
Then $\abs{(s_n+t_n)-(s+t)}\leq\abs{s_n-s}+\abs{t_n-t}<\varepsilon$.
\item Scalar multiplication uses $\abs{cs_n-cs}=\abs{c}\abs{s_n-s}$; translation is (i) with constant sequence $c$.
\item Since $s_n\to s$, choose $N\geq2$ with $\abs{s_n-s}<1$ for $n\geq N$ and set $M\defeq\max\{\abs{s_1},\ldots,\abs{s_{N-1}},\abs{s}+1\}$.\\
Then $\abs{s_nt_n-st}\leq \abs{s_n}\abs{t_n-t}+\abs{t}\abs{s_n-s}\leq M\abs{t_n-t}+\abs{t}\abs{s_n-s}\to0$.
\item If $s\neq0$, then for all sufficiently large $n$, $\abs{s_n}>\abs{s}/2$. Hence $\abs{\frac1{s_n}-\frac1s}=\frac{\abs{s_n-s}}{\abs{s_n}\abs{s}}\leq \frac{2}{\abs{s}^2}\abs{s_n-s}\to0$.
\end{enumerate}
\hlt{Structural Commentary}: Boundedness controls products; nonzero limits control reciprocals.\\
\hlt{Key Insight}: Product and reciprocal rules are estimates after the easy sum rule.
\end{proof}
\end{theorem}

\begin{theorem} \label{thm:rk-sequence-coordinate-convergence} \hlt{Coordinatewise Convergence in $\R^k$}\\
Let $\mathbf{x}_n=(\alpha_{j,n})_{j=1}^k$ and $\mathbf{x}=(\alpha_j)_{j=1}^k$ in $\R^k$. Then $\mathbf{x}_n\to\mathbf{x}$ $\Longleftrightarrow$ $\alpha_{j,n}\to\alpha_j$ $(1\leq j\leq k)$.\\
Moreover, if $\mathbf{x}_n\to\mathbf{x}$, $\mathbf{y}_n\to\mathbf{y}$, and $\beta_n\to\beta$ in $\R$, then $\mathbf{x}_n+\mathbf{y}_n\to\mathbf{x}+\mathbf{y}$; $\mathbf{x}_n\cdot\mathbf{y}_n\to\mathbf{x}\cdot\mathbf{y}$; $\beta_n\mathbf{x}_n\to\beta\mathbf{x}$.
\begin{proof}
\begin{enumerate}[label=\roman*.]
\setlength{\itemsep}{0pt}
\item[$(\Rightarrow)$] Since $\abs{\alpha_{j,n}-\alpha_j}\leq\norm{\mathbf{x}_n-\mathbf{x}}_2$ for each $j$, each coordinate converges.
\item[$(\Leftarrow)$] Given $\varepsilon>0$, choose $N$ so that $\abs{\alpha_{j,n}-\alpha_j}<\varepsilon/\sqrt{k}$ for all $j$ and $n\geq N$. Then $\norm{\mathbf{x}_n-\mathbf{x}}_2<\varepsilon$.
\end{enumerate}
The algebraic conclusions follow coordinatewise from Theorem \ref{thm:complex-sequence-limit-algebra}.\\
\hlt{Structural Commentary}: Euclidean convergence is exactly simultaneous coordinate convergence.\\
\hlt{Key Insight}: The norm packages finitely many coordinate estimates into one estimate.
\end{proof}
\end{theorem}

\begin{remark} \label{rem:coordinate-metric-convergence-equivalence} \hlt{Coordinate Metrics Give the Same Limits}\\
For $k\in\N$ and $x,y\in\R^k$, set $d_1(x,y)=\sum\limits_{j=1}^k\abs{x_j-y_j}$, $d_2(x,y)=\big(\sum\limits_{j=1}^k\abs{x_j-y_j}^2\big)^{1/2}$,  $d_\infty(x,y)=\max\limits_{1\leq j\leq k}\abs{x_j-y_j}$.\\
Since $d_\infty\leq d_2\leq d_1\leq k d_\infty$, these have same convergent sequences, equivalently same coordinatewise criterion.
\end{remark}

\begin{definition} \label{def:subsequence-and-subsequential-limit} \hlt{Subsequences and Subsequential Limits}\\
Let $(p_n)$ be a sequence. If $n_1<n_2<n_3<\cdots$ is a strictly increasing sequence of positive integers, then $(p_{n_k})$ is a \hlt{subsequence} of $(p_n)$. If $p_{n_k}\to p$, then $p$ is a \hlt{subsequential limit} of $(p_n)$.
\end{definition}

\begin{remark} \label{rem:subsequence-selection-map} \hlt{Subsequences as Selection Maps}\\
Equivalently, a subsequence is the composite of the original sequence, viewed as a function on $\N$ as in Definition \ref{def:sequence-as-enumeration}, with a strictly increasing map $\sigma:\N\to\N$. In this notation the subsequence is $k\mapsto p_{\sigma(k)}$.\\
\hlt{Key Insight}: A subsequence keeps the original order but is allowed to skip terms.
\end{remark}

\begin{remark} \label{rem:subsequences-of-convergent-sequences} \hlt{Subsequences of Convergent Sequences}\\
If $p_n\to p$, then every subsequence of $(p_n)$ also converges to $p$. Conversely, if every subsequence of $(p_n)$ converges to $p$, then $(p_n)$ itself converges to $p$, since $(p_n)$ is one of its subsequences.
\end{remark}

\begin{definition} \label{def:monotone-sequence} \hlt{Monotone Real Sequence}\\
A real sequence $(s_n)$ is \hlt{monotonically increasing} if $s_n\leq s_{n+1}$ for all $n$, and \hlt{monotonically decreasing} if $s_n\geq s_{n+1}$ for all $n$. A \hlt{monotonic} sequence is either increasing or decreasing.
\end{definition}

\begin{definition} \label{def:dominant-index} \hlt{Dominant Index}\\
For a real sequence $(s_n)$, an index $n$ is \hlt{dominant} if $s_n\geq s_m$ for every $m>n$.
\end{definition}

\begin{theorem} \label{thm:monotone-subsequence-theorem} \hlt{Monotone Subsequence Theorem}\\
Every real sequence has a monotone subsequence.
\begin{proof}
If the dominant indices $n_1<n_2<\cdots$ are infinite, then $s_{n_1}\geq s_{n_2}\geq\cdots$, giving a decreasing subsequence.\\
Otherwise choose $N$ beyond all dominant indices and set $n_1>N$.\\
Given $n_k>N$, nondominance gives $n_{k+1}>n_k$ with $s_{n_{k+1}}>s_{n_k}$; hence $(s_{n_k})$ is increasing.\\
\hlt{Structural Commentary}: Dominant indices give decreasing case; nondominance recursively gives increasing case.\\
\hlt{Key Insight}: The proof organizes arbitrary real behavior into one ordered thread.
\end{proof}
\end{theorem}

\begin{figure}[H]
  \centering
  \begin{tikzpicture}[
  >=Stealth,
  x=0.9cm,
  y=0.9cm,
  every node/.style={font=\small},
  axis/.style={<->, gray!70, line width=0.5pt},
  point/.style={circle, fill=black, inner sep=1.35pt},
  snpoint/.style={circle, fill=blue!70!black, inner sep=1.45pt},
  supoint/.style={circle, fill=green!50!black, inner sep=1.45pt}
]
  \draw[axis] (0.5,0) -- (7.55,0);

  \coordinate (L) at (4.45,0);
  \coordinate (S) at (5.65,0);
  \coordinate (R) at (6.85,0);
  \coordinate (SN) at (4.74,0);

  \node[gray!70] at (L) {$($};
  \node[gray!70] at (R) {$)$};

  \draw[green!50!black, dashed, line width=0.45pt] (S) -- ++(0,0.62);
  \node[below] at (L) {$s-\varepsilon$};
  \node[below] at (R) {$s+\varepsilon$};
  \node[green!50!black, anchor=west] at ($(S)+(0,0.78)$) {$s=\sup\{s_n:n\in\N\}$};

  \foreach \x/\lab in {1.75/s_1,2.50/s_2,3.25/s_3} {
    \node[point] at (\x,0) {};
    \node[above] at (\x,0.12) {$\lab$};
  }
  \node[above] at (3.90,0.12) {$\cdots$};

  \foreach \x in {5.08,5.30,5.44,5.53,5.59,5.63} {
    \node[point] at (\x,0) {};
  }
  \node[snpoint] at (SN) {};
  \node[supoint] at (S) {};

  \draw[blue!70!black, dashed, line width=0.45pt] (SN) -- ++(0,0.62);
  \node[blue!70!black] at ($(SN)+(0,0.78)$) {$s_N$};
\end{tikzpicture}

  \caption{For an increasing bounded sequence, $s=\sup\{s_n\}$ is the limit.}
  \label{fig:monotone-convergence-supremum}
\end{figure}

\begin{theorem} \label{thm:monotone-convergence} \hlt{Monotone Convergence Theorem for Sequences}\\
A monotonic real sequence converges if and only if it is bounded.
\begin{proof}
\begin{enumerate}[label=\roman*.]
\setlength{\itemsep}{0pt}
\item[$(\Rightarrow)$] If $s_n\to s$, choose $N$ with $\abs{s_n-s}<1$ for $n\geq N$. The tail lies in $(s-1,s+1)$ and the initial set $\{s_1,\ldots,s_{N-1}\}$ is finite, so $(s_n)$ is bounded.
\item[$(\Leftarrow)$] Let $E\defeq\{s_n:n\in\N\}$. If $(s_n)$ is increasing, set $s\defeq\sup E$. Given $\varepsilon>0$, choose $N$ with $s-\varepsilon<s_N\leq s$; then $s-\varepsilon<s_n\leq s$ for $n\geq N$, so $s_n\to s$. If $(s_n)$ is decreasing, set $s\defeq\inf E$. Given $\varepsilon>0$, choose $N$ with $s\leq s_N<s+\varepsilon$; then $s\leq s_n\leq s_N<s+\varepsilon$ for $n\geq N$, so $s_n\to s$.
\end{enumerate}
\hlt{Structural Commentary}: Monotonicity traps bounded tails at $\sup E$ or $\inf E$.\\
\hlt{Key Insight}: The least upper bound is forced to be the limit.
\end{proof}
\end{theorem}

\begin{theorem} \label{thm:bolzano-weierstrass-real-sequence} \hlt{Bolzano--Weierstrass for Real Sequences}\\
Every bounded real sequence has a convergent subsequence.
\begin{proof}
By Theorem \ref{thm:monotone-subsequence-theorem}, choose a monotone subsequence. It is bounded because the original sequence is bounded, so Theorem \ref{thm:monotone-convergence} gives convergence.\\
\hlt{Structural Commentary}: Boundedness prevents escape; monotone subsequences force one ordered thread.\\
\hlt{Key Insight}: Monotone subsequence plus monotone convergence is the real-line compactness mechanism.
\end{proof}
\end{theorem}

\begin{definition} \label{def:real-cauchy-sequence} \hlt{Cauchy Real Sequence}\\
A real sequence $(s_n)$ is \hlt{Cauchy} if, for every $\varepsilon>0$, there is $N\in\N$ such that $m,n\geq N$ implies $\abs{s_n-s_m}<\varepsilon$.
\end{definition}

\begin{lemma} \label{lem:real-cauchy-sequence-with-convergent-subsequence} \hlt{Cauchy Plus One Convergent Subsequence}\\
If a real Cauchy sequence has a subsequence converging to $s\in\R$, then the whole sequence converges to $s$.
\begin{proof}
Given $\varepsilon>0$, choose $N_0$ so that $m,n\geq N_0$ implies $\abs{s_m-s_n}<\varepsilon/2$. Choose $j$ with $n_j\geq N_0$ and $\abs{s_{n_j}-s}<\varepsilon/2$. Then for every $n\geq N_0$, $\abs{s_n-s}\leq\abs{s_n-s_{n_j}}+\abs{s_{n_j}-s}<\varepsilon$.\\
\hlt{Structural Commentary}: Cauchy tails follow any convergent thread.\\
\hlt{Key Insight}: Mutual tail closeness turns one subsequential limit into the full limit.
\end{proof}
\end{lemma}

\begin{theorem} \label{thm:real-cauchy-convergence-criterion} \hlt{Real Cauchy Criterion}\\
A real sequence converges in $\R$ if and only if it is Cauchy.
\begin{proof}
If $s_n\to s$, choose $N$ such that $\abs{s_n-s}<\varepsilon/2$ for $n\geq N$; then $\abs{s_n-s_m}<\varepsilon$ for $m,n\geq N$. Conversely, suppose $(s_n)$ is Cauchy. Choose $N$ so that $m,n\geq N$ implies $\abs{s_m-s_n}<1$. Then the tail lies in $(s_N-1,s_N+1)$, and the finitely many initial terms are bounded, so $(s_n)$ is bounded. By Theorem \ref{thm:bolzano-weierstrass-real-sequence}, some subsequence converges; Lemma \ref{lem:real-cauchy-sequence-with-convergent-subsequence} makes the whole sequence converge.\\
\hlt{Structural Commentary}: Completeness of $\R$ means Cauchy behavior has a real endpoint.\\
\hlt{Key Insight}: Cauchy gives boundedness, Bolzano--Weierstrass gives a subsequential endpoint, and the Cauchy property pulls the whole sequence there.
\end{proof}
\end{theorem}

\begin{definition} \label{def:sequence-diverges-to-infinity} \hlt{Divergence to $\pm\infty$}\\
For a real sequence $(s_n)$,
\begin{align}
s_n\to+\infty &\iff \forall M\in\R,\ \exists N\in\N:\ n\geq N\Rightarrow s_n\geq M, \nonumber\\
s_n\to-\infty &\iff \forall M\in\R,\ \exists N\in\N:\ n\geq N\Rightarrow s_n\leq M. \nonumber
\end{align}
This extends $\to$ notation for divergent real sequences; Definition \ref{def:real-complex-sequence-convergence} is unchanged.
\end{definition}

\begin{proposition} \label{prop:unbounded-real-sequence-has-infinite-subsequence} \hlt{Unbounded Real Sequences Have Infinite Subsequences}\\
A real sequence is unbounded iff some subsequence tends to $+\infty$ or $-\infty$.
\begin{proof}
Set $n_0=0$. If $(s_n)$ is unbounded above, every tail is unbounded above, so choose $n_k>n_{k-1}$ with $s_{n_k}\geq k$; then $s_{n_k}\to+\infty$. The unbounded-below case similarly gives $s_{n_k}\leq-k$, hence $s_{n_k}\to-\infty$.\\
Conversely, a subsequence tending to $\pm\infty$ is unbounded, so $(s_n)$ is unbounded.\\
\hlt{Structural Commentary}: Unbounded tails supply the thresholds $\pm k$.\\
\hlt{Key Insight}: Large late terms form a divergent subsequence.
\end{proof}
\end{proposition}

\begin{definition} \label{def:upper-and-lower-limits} \hlt{Upper and Lower Limits}\\
For real sequence $(s_n)$, let $E\defeq\{x\in\overline{\R}:\exists\, n_1<n_2<\cdots,\ s_{n_k}\to x\}$, allowing $x=\pm\infty$.\\
Set $s^*\defeq\sup E$ and $s_*\defeq\inf E$; equivalently, $\limsup_{n\to\infty}s_n=s^*$ and $\liminf_{n\to\infty}s_n=s_*$.
\end{definition}

\begin{theorem} \label{thm:limsup-characterization} \hlt{Characterization of the Upper Limit}\\
Let $(s_n)$ be a real sequence, and let $E,s^*,s_*$ be as in Definition \ref{def:upper-and-lower-limits}. Then:
\begin{enumerate}[label=(\roman*)]
\setlength{\itemsep}{0pt}
\item $s^*\in E$.
\item If $x>s^*$, then $s_n<x$ eventually.
\item These two properties characterize $s^*$ uniquely.
\end{enumerate}
Dually, $s_*$ is uniquely characterized by $s_*\in E$ and $x<s_*$ implies $s_n>x$ eventually.
\begin{proof}
\begin{enumerate}[label=(\roman*)]
\setlength{\itemsep}{0pt}
\item Since $s^*=\sup E$: if $s^*=+\infty$, use Proposition \ref{prop:unbounded-real-sequence-has-infinite-subsequence}; if $s^*\in\R$, use Theorem \ref{thm:subsequential-limits-closed}. If $s^*=-\infty$ and $s_n\nrightarrow-\infty$, infinitely many terms are bounded below; they are either unbounded above or have a bounded subsequence, giving an element of $E$ above $-\infty$, a contradiction.
\item If $s_n\geq x>s^*$ infinitely often, that subsequence is unbounded above or has a finite subsequential limit $y\geq x$; both contradict $s^*=\sup E<x$. Hence $s_n<x$ eventually.
\item If $p<q$ both satisfy (i)--(ii), choose $p<x<q$. Then $s_n<x$ eventually, contradicting $q\in E$.
\end{enumerate}
Apply the same argument to $(-s_n)$ for $s_*$.\\
\hlt{Structural Commentary}: Limsup is the maximal subsequential limit and eventual ceiling.\\
\hlt{Key Insight}: Limsup combines subsequences with eventual upper bounds.
\end{proof}
\end{theorem}

\begin{theorem} \label{thm:convergence-iff-limsup-liminf-equal} \hlt{Convergence and Equal Upper/Lower Limits}\\
A real sequence $(s_n)$ converges to a finite number $s\in\R$ if and only if $\limsup_{n\to\infty}s_n=\liminf_{n\to\infty}s_n=s$.
\begin{proof}
\begin{enumerate}[label=\roman*.]
\setlength{\itemsep}{0pt}
\item[$(\Rightarrow)$] If $s_n\to s$, every subsequence converges to $s$, and the tail is bounded, so no subsequence tends to $\pm\infty$. Thus $E=\{s\}$, hence $\limsup s_n=\liminf s_n=s$.
\item[$(\Leftarrow)$] Suppose $\limsup s_n=\liminf s_n=s$. Given $\varepsilon>0$, Theorem \ref{thm:limsup-characterization} gives $s_n<s+\varepsilon$ and $s_n>s-\varepsilon$ eventually. Hence $\abs{s_n-s}<\varepsilon$ eventually, so $s_n\to s$.
\end{enumerate}
\hlt{Structural Commentary}: Equality of the two extremal subsequential limits traps the whole tail.\\
\hlt{Key Insight}: Limsup supplies an eventual ceiling; liminf supplies an eventual floor.
\end{proof}
\end{theorem}

\begin{theorem} \label{thm:limsup-liminf-order-comparison} \hlt{Order Comparison for Upper and Lower Limits}\\
If $s_n\leq t_n$ for all $n\geq N$, where $N$ is fixed, then:
\begin{enumerate}[label=(\roman*)]
\setlength{\itemsep}{0pt}
\item $\limsup\limits_{n\to\infty}s_n\leq \limsup\limits_{n\to\infty}t_n$.
\item $\liminf\limits_{n\to\infty}s_n\leq \liminf\limits_{n\to\infty}t_n$.
\end{enumerate}
\begin{proof}
Discard finitely many terms, so $s_n\leq t_n$ for all $n$.
\begin{enumerate}[label=(\roman*)]
\setlength{\itemsep}{0pt}
\item Put $a=\limsup s_n$ and $b=\limsup t_n$. If $x<a$, then $s_n>x$ infinitely often, hence $t_n>x$ infinitely often; Theorem \ref{thm:limsup-characterization} gives $x\leq b$. Thus $a\leq b$.
\item Apply (i) to $-t_n\leq -s_n$ to get $\limsup(-t_n)\leq\limsup(-s_n)$, equivalently $\liminf s_n\leq\liminf t_n$.
\end{enumerate}
\hlt{Structural Commentary}: Eventual order is inherited by the extremal subsequential limits.\\
\hlt{Key Insight}: Limsup and liminf ignore finitely many initial terms.
\end{proof}
\end{theorem}

\begin{theorem} \label{thm:pinching-theorem-sequences} \hlt{Pinching Theorem for Real Sequences}\\
Let $(a_n)$, $(s_n)$, and $(b_n)$ be real sequences with $a_n\leq s_n\leq b_n$ for all sufficiently large $n$.\\
If $a_n\to s$ and $b_n\to s$, then $s_n\to s$.
\begin{proof}
Let $\varepsilon>0$. Choose $N_0$ s.t. $a_n\leq s_n\leq b_n$ for $n\geq N_0$; $N_1$ s.t. $s-\varepsilon<a_n$ for $n\geq N_1$; and $N_2$ s.t. $b_n<s+\varepsilon$ for $n\geq N_2$. If $N=\max\{N_0,N_1,N_2\}$, then $s-\varepsilon<s_n<s+\varepsilon$ for all $n\geq N$. Hence $s_n\to s$.\\
\hlt{Structural Commentary}: Two convergent bounds squeeze the middle sequence to the same limit.\\
\hlt{Key Insight}: Once both walls are inside the $\varepsilon$-neighborhood, the middle term has no room to escape.
\end{proof}
\end{theorem}

\begin{figure}[H]
  \centering
  \begin{tikzpicture}[
  >=Stealth,
  x=0.9cm,
  y=0.9cm,
  every node/.style={font=\small},
  axis/.style={->, gray!70, line width=0.55pt},
  upper/.style={blue!70!black, line width=0.9pt},
  lower/.style={orange!85!black, line width=0.9pt},
  upperdot/.style={circle, fill=blue!70!black, inner sep=1.35pt},
  lowerdot/.style={circle, fill=orange!85!black, inner sep=1.35pt},
  middledot/.style={circle, draw=green!45!black, fill=white, line width=0.65pt, inner sep=1.45pt}
]
  \draw[axis] (0.35,0) -- (7.55,0) node[right] {$n$};
  \draw[axis] (0.7,-1.45) -- (0.7,1.55);
  \node[left] at (0.65,0) {$s$};

  \coordinate (u1) at (1,1.18);
  \coordinate (u2) at (2,0.72);
  \coordinate (u3) at (3,0.45);
  \coordinate (u4) at (4,0.28);
  \coordinate (u5) at (5,0.18);
  \coordinate (u6) at (6,0.10);
  \coordinate (u7) at (7,0.05);

  \draw[upper]
    plot[smooth] coordinates {(u1) (u2) (u3) (u4) (u5) (u6) (u7)};

  \coordinate (l1) at (1,-1.08);
  \coordinate (l2) at (2,-0.60);
  \coordinate (l3) at (3,-0.34);
  \coordinate (l4) at (4,-0.20);
  \coordinate (l5) at (5,-0.11);
  \coordinate (l6) at (6,-0.06);
  \coordinate (l7) at (7,-0.03);

  \draw[lower]
    plot[smooth] coordinates {(l1) (l2) (l3) (l4) (l5) (l6) (l7)};

  \foreach \p in {u1,u2,u3,u4,u5,u6,u7} {
    \node[upperdot] at (\p) {};
  }
  \foreach \p in {l1,l2,l3,l4,l5,l6,l7} {
    \node[lowerdot] at (\p) {};
  }

  \foreach \x/\ym in {
    1/-0.42,
    2/0.16,
    3/0.02,
    4/0.12,
    5/-0.02,
    6/0.03,
    7/0.01
  } {
    \node[middledot] at (\x,\ym) {};
  }

  \node[blue!70!black, anchor=west] at (1.08,1.20) {$(b_n)$};
  \node[green!45!black, anchor=west] at (1.08,-0.42) {$(s_n)$};
  \node[orange!85!black, anchor=west] at (1.08,-1.2) {$(a_n)$};
\end{tikzpicture}

  \caption{With $a_n\leq s_n\leq b_n$ and $a_n,b_n\to s$, the middle sequence also tends to $s$.}
  \label{fig:pinching-theorem-sequences}
\end{figure}

\begin{theorem} \label{thm:special-sequence-limits} \hlt{Special Sequence Limits}
\begin{enumerate}[label=(\roman*)]
\setlength{\itemsep}{0pt}
\item If $p>0$, then $\lim\limits_{n\to\infty}n^{-p}=0$.
\item If $p>0$, then $\lim\limits_{n\to\infty}\sqrt[n]{p}=1$.
\item $\lim\limits_{n\to\infty}\sqrt[n]{n}=1$.
\item If $p>0$ and $\alpha\in\R$, then $\lim\limits_{n\to\infty}\frac{n^\alpha}{(1+p)^n}=0$.
\item If $\abs{x}<1$, then $\lim\limits_{n\to\infty}x^n=0$.
\end{enumerate}
\begin{proof}
\begin{enumerate}[label=(\roman*)]
\setlength{\itemsep}{0pt}
\item Given $\varepsilon>0$, choose $N>(1/\varepsilon)^{1/p}$ by Theorem \ref{thm:archimedean-and-rational-density}. Then $n\geq N$ gives $n^{-p}<\varepsilon$.
\item If $p>1$, set $x_n\defeq\sqrt[n]{p}-1$. Then $x_n>0$ and $1+nx_n\leq(1+x_n)^n=p$, so $0<x_n\leq(p-1)/n\to0$.\\
By Theorem \ref{thm:pinching-theorem-sequences}, $x_n\to0$. If $p=1$, then $\sqrt[n]{p}=1$ for every $n$.\\
If $0<p<1$, then $1/p>1$, so $\sqrt[n]{1/p}\to1$ by the previous case.\\
Since $\sqrt[n]{p}=1/\sqrt[n]{1/p}$ and limits preserve reciprocals away from $0$, $\sqrt[n]{p}\to1$.
\item Put $x_n\defeq\sqrt[n]{n}-1$. Then $x_n\geq0$, and for $n\geq2$ the binomial theorem gives $n=(1+x_n)^n\geq \frac{n(n-1)}{2}x_n^2$.\\
Thus $0\leq x_n\leq\sqrt{2/(n-1)}$, so Theorem \ref{thm:pinching-theorem-sequences} gives $x_n\to0$.
\item Choose an integer $k>\alpha$ with $k>0$. For $n>2k$, $(1+p)^n\geq \binom{n}{k}p^k=\frac{n(n-1)\cdots(n-k+1)}{k!}p^k>\frac{n^k p^k}{2^k k!}$.\\
Hence $0<\frac{n^\alpha}{(1+p)^n}<\frac{2^kk!}{p^k}n^{\alpha-k}$, and $n^{\alpha-k}\to0$ by (i). Apply Theorem \ref{thm:pinching-theorem-sequences}.
\item If $x\neq0$, write $\abs{x}=(1+p)^{-1}$ with $p>0$ and apply (iv) with $\alpha=0$ to get $\abs{x}^n\to0$. Since $-\abs{x}^n\leq x^n\leq\abs{x}^n$, Theorem \ref{thm:pinching-theorem-sequences} gives $x^n\to0$; $x=0$ is immediate.
\end{enumerate}
\hlt{Structural Commentary}: Each limit is reduced to an elementary comparison sequence.\\
\hlt{Key Insight}: The binomial theorem supplies the right upper bounds for roots and exponential domination.
\end{proof}
\end{theorem}

\begin{method} \label{meth:sequence-proof-patterns}\label{meth:sequence-proof-patterns-source-todo} \hlt{Sequence Proof Patterns}
\begin{enumerate}[label=(\roman*)]
\setlength{\itemsep}{0pt}
\item \hlt{Convergence}: fix $\varepsilon>0$, then find $N$ so that all later terms lie in the relevant $\varepsilon$-ball.
\item \hlt{Uniqueness}: compare two proposed limits with $\varepsilon/2$ and use the triangle inequality.
\item \hlt{Subsequences}: choose $n_k\uparrow$ recursively with distances $<1/k$.
\item \hlt{Cauchy arguments}: estimate the diameter of tails rather than guessing the limit first.
\item \hlt{Monotone sequences}: use $\sup$ or $\inf$ and then prove eventual trapping.
\end{enumerate}
\hlt{Key Insight}: Most sequence proofs are bookkeeping for quantifiers, indices, and one controlling inequality.
\end{method}
